\documentclass[11pt]{article}

% ---- ICLR-like page setup (self-contained; no external style file) ----
\usepackage[letterpaper,top=1in,bottom=1in,left=1.5in,right=1.5in]{geometry}
\usepackage{times}
\usepackage[T1]{fontenc}
\usepackage[utf8]{inputenc}
\usepackage{amsmath,amssymb,amsthm}
\usepackage{booktabs}
\usepackage{graphicx}
\usepackage{caption}
\usepackage{subcaption}
\usepackage{url}
\usepackage{hyperref}
\usepackage{xcolor}
\usepackage{microtype}
\usepackage{enumitem}
\usepackage{multirow}
\usepackage[authoryear,round]{natbib}
\hypersetup{colorlinks=true,linkcolor=blue!60!black,citecolor=blue!60!black,urlcolor=blue!60!black}
\setlength{\parskip}{3pt}
\newcommand{\Rb}{\mathbb{R}}
\newcommand{\E}{\mathbb{E}}
\newcommand{\xh}{\hat{x}}
\newcommand{\todo}[1]{\textcolor{red}{[#1]}}
\newtheorem{proposition}{Proposition}
\newcommand{\figdir}{../runs/wide_rp_300/figures/}
\newcommand{\diagram}[2]{\IfFileExists{figures/#1.pdf}{\includegraphics[width=#2]{figures/#1.pdf}}{\fbox{\texttt{figures/#1.pdf}}}}
\newcommand{\figorplaceholder}[2]{%
  \IfFileExists{\figdir#1}{\includegraphics[width=#2]{\figdir#1}}{\fbox{\parbox{#2}{\centering\vspace{2em}\texttt{#1}\\(produced by \texttt{just paper})\vspace{2em}}}}}

\title{\textbf{Layout-Conditioned Flow Matching with Mask-Guided Region Pooling\\for Anime Face Generation:\\A Complete Record of Methods and Decisions}}
\author{tinyflow project \\ \texttt{feat/pixel-unet}, September 2026 (updated 2026-09-20)}
\date{}

\begin{document}
\maketitle

\begin{abstract}
We document, end to end, the construction of a $64\times64$ anime-face generator based on
flow matching with x-prediction in pixel space, trained from scratch on 21{,}551 images
with one RTX~4090. The record covers the data pipeline, the network to the level of
every layer's normalisation and initialisation, the objective, optimiser, sampler and
evaluation protocol, and---at equal weight---every design decision, including the ones
that failed. The central negative result is a derivation, confirmed by ablation, that
no auxiliary loss on the predicted clean image against the target (edges, lines,
palettes, mask-weighted variants, per-pixel re-weightings) can improve a flow-matching
model: at the objective's optimum such penalties are almost entirely irreducible and
only bias the network toward hedged predictions. The central positive result is a
structural change: conditioning on a three-channel semantic layout (face, eyes, mouth
hulls from a landmark detector), sampled at test time from a Gaussian-mixture prior over
a point-distribution model of the landmarks, together with \emph{mask-guided region
pooling}---a zero-initialised layer that pools decoder features inside each semantic
region and broadcasts a projection back into that region only. It reduces the rate of
mismatched iris colours from $35\%$ to the data's own $4$--$5\%$ at a cost of $3.5$ FID
against a schedule-matched unconditioned control. We also report where the model fails
(coverage: recall $0.15$ against a real-vs-real ceiling of $0.77$), the learning-rate
instability that shaped the final training schedule, and three sampling-time methods
that need no retraining: edge guidance ($-2$ to $-8$ FID), re-noise-and-re-solve
refinement ($-3$), and \emph{autoguidance} with a 40-epoch checkpoint of the same
recipe as the guide ($34.6\to21.2$; $33.2\to18.1$ on the best checkpoint), which
takes the layout-conditioned model $11$ FID past the unconditioned control. The
recipe transfers to CelebAMask-HQ at $64\times64$ with real label maps as layouts
(FID $27.3$, $18.6$ with autoguidance).
\end{abstract}

% =====================================================================
\section{Introduction}

This paper is a methods record rather than a claim of state of the art. Its purpose is
to let a reader reproduce every number, understand why each component is what it is,
and avoid repeating experiments that were tried and closed. It accompanies the
\texttt{tinyflow} repository; where a symbol is defined in code, the file and function
are named.

The starting point was a working pixel-space flow-matching model that shipped with two
domain-motivated auxiliary losses (a Sobel gradient-difference loss and a dark-line
``ink'' loss), the intuition being that anime is line art over flat fills. The campaign
asked whether such dataset-specific inductive biases help. Answering that required
(i)~a theory of what an auxiliary loss can do at the flow-matching optimum
(\S\ref{sec:theory}), (ii)~an offline test-bed that does not need training
(\S\ref{sec:offline}), (iii)~paired ablations with a measured noise floor
(\S\ref{sec:ablations}), and (iv)~a diagnosis of where the best model actually fails
(\S\ref{sec:bottleneck}). The diagnosis pointed at two distinct problems---a
long-range consistency failure (mismatched iris colours) and a coverage failure---and the
first of these is solved by the architecture of \S\ref{sec:layout}.

Contributions, in the order they were established:
\begin{enumerate}[leftmargin=1.5em,itemsep=1pt]
\item A derivation (\S\ref{sec:theory}) that training-time losses on $\xh$ against $x_1$
      are irreducible at the conditional-mean optimum, and that $x_1$-dependent per-pixel
      weights tilt the learned field; both confirmed by ablation (\S\ref{sec:ablations}).
\item A measurement protocol (\S\ref{sec:protocol}): FID noise, loss as the sensitive
      statistic for architecture changes, a validated iris-mismatch metric, and a
      16-step sampler for evaluation.
\item A layout prior over detector landmarks (\S\ref{sec:prior}) and a conditioned model
      with mask-guided region pooling (\S\ref{sec:regionpool}), the first structural
      change with an unambiguous, reproduced win.
\item A bottleneck analysis (\S\ref{sec:bottleneck}) locating the remaining gap in
      coverage and in small low-variance features (mouth, nose).
\end{enumerate}

% =====================================================================
\section{Data}\label{sec:data}

\paragraph{Images.} 21{,}551 anime face crops at $64\times64\times3$
(\texttt{data/animefaces.py}). PNGs are loaded once, converted to
\texttt{float32}, mapped to $[-1,1]$ by $x/127.5-1$ and cached
(\texttt{.preprocessed/anime\_faces.npy}). All public boundaries use $(H,W,C)$ layout;
only the U-Net's interior is channels-first. Images are denormalised for display with
$\mathrm{clip}((x+1)\cdot127.5, 0, 255)$; clipping is mandatory because the network has
no bounded output activation (\S\ref{sec:unet}).

\paragraph{Augmentation.} A Grain pipeline applies, per image, on worker threads:
a horizontal flip with probability $0.5$ (applied jointly to the image and its layout
mask), then colour jitter with brightness shift $b\sim U(-0.2,0.2)$, contrast
$c\sim U(0.8,1.2)$, saturation $s\sim U(0.8,1.2)$, as
$x \leftarrow x + b$; $x \leftarrow \bar{x} + c\,(x-\bar{x})$ with $\bar{x}$ the image
mean; $x \leftarrow g + s\,(x - g)$ with $g$ the per-pixel channel mean; then clipped to
$[-1,1]$. Masks pass through the jitter untouched. Batch size 128; batches straddle
epoch boundaries (ceiling division), so an epoch is $\lceil N/128\rceil$ batches.

\paragraph{Landmarks and masks.}\label{sec:masks} Each image is upsampled to $256\times256$
(bicubic) and passed to \texttt{hysts/anime-face-detector} (YOLOv3 face detector,
HRNet 28-keypoint regressor) with the whole image as the face box. The 28 points are
fixed in number and order: 0--4 chin contour left to right; 5--7 and 8--10 the two
brows left to right; 11--16 and 17--22 the two eyes, each as the upper lid left to
right (3 points) then the lower lid left to right (3 points); 23 nose; 24--27 mouth
(left corner, top, right corner, bottom). Coordinates are stored in the image's
$[-1,1]^2$ frame with the pixel-centre convention $((p+0.5)/256)\cdot 2-1$; points the
detector places outside the crop are kept (they occur for tight crops). Per-point
confidences $\in[0,1]$ are stored alongside
(\texttt{experiments/extract\_landmarks.py}).

The three-channel mask used for conditioning is rasterised from the landmarks: channel 0
the convex hull of points 0--10 (face), channel 1 the union of the hulls of 11--16 and
17--22 (eyes), channel 2 the hull of 24--27 (mouth). Hulls are filled at $256\times256$
and area-averaged to $64\times64$, giving soft, anti-aliased edges; values are stored as
\texttt{uint8} and used in $[0,1]$. The eye channel covers $\approx4.8\%$ of pixels.
\texttt{data/layouts.py:rasterize} reproduces the stored masks with IoU $0.99$ (face),
$0.95$ (eyes), $0.64$ (mouth; a few-pixel region, quantisation).

\paragraph{Curation.} Sorting images by mean landmark confidence and inspecting bands
(\texttt{runs/ablation/curation\_*.png}): below $0.3$ the content is dominated by
non-faces (torsos, clothing, hands, duplicated crops); the $0.3$--$0.65$ band is
legitimate hard faces (closed eyes, glasses, masks, mascots) that a generator should
keep. Training therefore drops images with mean confidence $<0.3$ (317 images,
$1.47\%$; 21{,}234 remain, 166 batches per epoch). \textbf{FID reference statistics
always use the full set}, so scores stay comparable across runs. Measured effect
(\S\ref{sec:ablations}): neutral to slightly positive.

% =====================================================================
\section{Model}\label{sec:model}

\subsection{Flow matching with x-prediction}\label{sec:fm}

With $x_1\sim p_{\text{data}}$, $x_0\sim\mathcal N(0,I)$ and a time $t\in[0,1]$, the
interpolant is $x_t=t\,x_1+(1-t)\,x_0$ and the conditional velocity is $x_1-x_0$. The
network predicts the \emph{clean image} $\xh=f_\theta(x_t,t)$ (x-prediction), and the
velocity is derived wherever it is needed as
\begin{equation}
u_\theta(x_t,t)=\frac{\xh-x_t}{\max(1-t,\ \epsilon_d)},\qquad \epsilon_d=0.05 .
\label{eq:vel}
\end{equation}
The training loss is the plain conditional flow-matching loss in velocity space,
\begin{equation}
\mathcal L(\theta)=\E_{t,x_0,x_1}\ \tfrac12\bigl\|u_\theta(x_t,t)-(x_1-x_0)\bigr\|^2
 =\E\ \tfrac12\frac{\|\xh-x_1\|^2}{\max(1-t,\epsilon_d)^2},
\label{eq:loss}
\end{equation}
i.e.\ a per-pixel squared error on $\xh$ with weight $w(t)=1/\max(1-t,\epsilon_d)^2$
($\approx1.2$ at $t=0.1$, $4$ at $0.5$, $100$ at $0.9$, $400$ at the floor). The floor
bounds the weight so that a single sample near $t=1$ cannot own a batch's gradient;
this is the JiT recipe \citep{li2026jit}. $t$ is drawn logit-normal,
$t=\sigma(\mu+\sigma_t z)$, $z\sim\mathcal N(0,1)$, $\mu=-0.8$, $\sigma_t=1$, which
puts more mass at low $t$ (high noise) and thins the tail at $t\to1$ where $w(t)$ is
large. The population minimiser of \eqref{eq:loss} at every $(x_t,t)$ is the
conditional mean $x^*(t)=\E[x_1\mid x_t]$ (a fact used throughout \S\ref{sec:theory}).

\textbf{The loss has no auxiliary terms.} This was not the starting point; it is the
conclusion of \S\ref{sec:theory}--\ref{sec:ablations}.

\subsection{Conditioning token}\label{sec:cond}

For a conditioned model with $K=3$ mask channels, the network input is the
concatenation along channels
\begin{equation}
x_{\text{in}}=[\,x_t;\ m\odot\kappa;\ \kappa\,]\in\Rb^{H\times W\times(3+K+1)},
\end{equation}
where $m\in[0,1]^{H\times W\times K}$ is the layout mask and $\kappa\in\{0,1\}$ a
per-sample indicator broadcast to a channel (\texttt{imagefm.cond\_token}). The
\emph{null token} is $\kappa=0$ with zero mask channels; the indicator disambiguates
``unconditioned'' from ``a face with no parts''. Classifier-free training
\citep{ho2022cfg} replaces the mask by the null token with probability $p_{\text{drop}}$
(\texttt{cond\_dropout}); the final recipe uses $p_{\text{drop}}=0$ (a pure conditional
model) because dropout cost $\approx7$ FID at the screening horizon
(\S\ref{sec:condres}). The target is then $\E[x_1\mid x_t,m]$: the layout is given and
the network renders. Sampling never uses a real image's mask; layouts come from the
prior of \S\ref{sec:prior}.

\subsection{Sampler}

The probability-flow ODE $\dot x=u_\theta(x,t)$ is integrated from $t=0$ to
$t_1=1-10^{-3}$ with \texttt{diffrax}'s Dopri5 on a fixed grid of $n$ steps
($\Delta t=t_1/n$, no adaptive control), $n=64$ for reported samples and $n=16$ for
FID during training and screening (\S\ref{sec:protocol}). One solver
(\texttt{imagefm.\_solve}) serves training-time samples, FID and figures. For a
conditioned model the conditioning token is held fixed along the trajectory.

\subsection{Network}\label{sec:unet}

\begin{figure}[t]
\centering
\diagram{architecture}{\linewidth}
\caption{The velocity network. The image $x_t$ is the only input to the U-Net;
the layout masks $m$ enter solely through the region-pooling layers after the
decoder levels producing $16^2$ and $32^2$. Four encoder levels (ResBlock doubling
the channels, then a $2\times2$ stride-2 depthwise convolution), a bare
$4^2\times1024$ tensor with no mid block, four decoder levels (bilinear
$\uparrow2$ and a $3\times3$ convolution, concatenation of the skip, ResBlock
halving the channels). Skips run from the \emph{input} of each encoder level to the
decoder level after its upsample (Fig.~\ref{fig:skips}). One shared time embedding
$e(t)$ feeds a separate adaLN-Zero projection in every ResBlock; the output is
$\hat x$ with no activation, and the velocity is derived from it.}
\label{fig:architecture}
\end{figure}

The velocity network is a four-level U-Net (\texttt{models/unet.py}), channels-first
internally, $36.7$\,M parameters at \texttt{base\_channels}$=64$ ($9.38$\,M at $32$).
\begin{itemize}[leftmargin=1.5em,itemsep=1pt]
\item \textbf{Input/output.} A $3\times3$ convolution (padding 1) maps the
      $3+K+1=7$ input channels to $64$; a $1\times1$ convolution maps the final $64$
      channels to the $3$ output channels. There is \textbf{no final activation}: the
      velocity derived from $\xh$ is unbounded regardless, and a bounded output would
      only add vanishing gradients.
\item \textbf{Time embedding.} $t$ is scaled by $10^3$ and embedded with 64
      frequencies $\omega_k$ geometrically spaced from $1$ to $10^{-4}$,
      $e(t)=[\sin(10^3 t\,\omega_k);\cos(10^3 t\,\omega_k)]\in\Rb^{128}$, then a
      two-layer MLP $128\to128\to128$ with SiLU. Without the $10^3$ scale no frequency
      completes a cycle on $[0,1]$ and the embedding is nearly constant in $t$.
\item \textbf{ResBlock} (in $c$, out $c'$): GroupNorm(8 groups, $\epsilon=10^{-5}$,
      statistics over the group's channels \emph{and} space, learned affine) $\to$ SiLU
      $\to$ $3\times3$ conv $\to$ GroupNorm (no affine) $\to$ adaptive modulation
      $h\odot(1+\gamma)+\beta$ with $[\gamma;\beta]=W_a\,\mathrm{SiLU}(e)+b_a$,
      $W_a\in\Rb^{2c'\times128}$ \textbf{zero-initialised} $\to$ SiLU $\to$ $3\times3$
      conv, plus a skip path ($1\times1$ conv if $c\neq c'$, identity otherwise). The
      zero initialisation follows adaLN-Zero \citep{peebles2023dit}, the FiLM analogue
      of the zero-$\gamma$ trick of \citet{wu2018gn}; the modulation has no output
      gate. GroupNorm \citep{wu2018gn} rather than InstanceNorm is deliberate: per-channel
      statistics would strip per-image brightness and colour cast, which a generator
      must keep.
\item \textbf{UBlock} at level $i$ with channels $c_i=64\cdot2^i$: \emph{down} =
      ResBlock($c_i\to2c_i$) $\to$ SiLU $\to$ depthwise $2\times2$ stride-2 conv;
      \emph{up} = bilinear $\times2$ $\to$ $3\times3$ conv ($2c_i\to c_i$) $\to$ SiLU
      $\to$ concatenate the skip ($c_i$ channels) $\to$ ResBlock($2c_i\to c_i$).
      Resolutions $64,32,16,8$ at the four levels and a $4\times4\times1024$
      bottleneck; the bottleneck is the last down-conv's output with no mid-block.
\item \textbf{Skip layout (a documented deviation).} The skip saved at level $i$ is the
      \emph{input} to that level's down block, so level $i{+}1$'s skip is level $i$'s
      encoder output after the strided conv---the standard pattern shifted down one
      level: the top decoder level receives \texttt{in\_conv} features rather than a
      ResBlock output, and each skip is half the standard width. Tested
      (\S\ref{sec:arch}): the standard layout is neutral at $9.4$\,M parameters and
      unstable at $37$\,M at the default learning rate; the deviation is kept.
\item \textbf{Region pooling} (\S\ref{sec:regionpool}) after the up blocks that
      produce $16\times16$ (256 channels) and $32\times32$ (128 channels).
\end{itemize}

\subsection{Mask-guided region pooling}\label{sec:regionpool}

\begin{figure}[t]
\centering
\diagram{regionpool}{\linewidth}
\caption{The region-pooling layer. Features are mean-pooled inside each soft region
(face, eyes, mouth, hair/background), projected by a per-region zero-initialised
$C\times C$ matrix and broadcast back into that region, as a residual. Both irises
are painted from the same pooled eye vector.}
\label{fig:regionpool}
\end{figure}

\begin{figure}[t]
\centering
\diagram{skips}{\linewidth}
\caption{Skip wiring. Standard U-Nets pass each encoder level's output to the
matching decoder level; this network passes the level's \emph{input} (the previous
level's output, half the width), i.e. the standard pattern shifted down one level.
Neutral at $9.4$\,M and the more stable choice at $37$\,M (\S\ref{sec:arch}).}
\label{fig:skips}
\end{figure}

Let $h\in\Rb^{C\times H\times W}$ be the decoder feature map after an up block and
$m_k\in[0,1]^{H\times W}$, $k=1..4$, the layout masks area-averaged to $H\times W$ for
the regions \emph{face}, \emph{eyes}, \emph{mouth} and \emph{hair/background}
$=1-\text{face}$ (read directly from the network's own input channels). The layer is
\begin{equation}
\bar e_k=\frac{\sum_p m_k(p)\,h(p)}{\sum_p m_k(p)+10^{-6}}\in\Rb^{C},\qquad
h'=h+\sum_{k=1}^{4} m_k\otimes\bigl(W_k\bar e_k+b_k\bigr),
\qquad W_k\in\Rb^{C\times C},\ b_k\in\Rb^{C},
\label{eq:rp}
\end{equation}
with $W_k=0,\ b_k=0$ at initialisation (the layer is the identity at step 0), all regions
pooled from the same input $h$ (order independent), implemented as three einsums.
Parameters: $4(C^2+C)$ per level, $0.33$\,M in total at width 64.

\emph{Interpretation.} In attention terms this is masked attention with a single fixed
query per region and uniform weights: the \emph{routing} (which positions belong
together) is supplied by the layout, and only the map on the aggregate is learned.
Both irises receive an identical additive term derived from both of them, so the two
local renderers start from the same offset and---with zero initialisation---can only be
moved together. The closest published relative is SEAN \citep{zhu2020sean}, which
broadcasts a per-region style code taken from an encoder of a reference image; here the
code is pooled from the network's own decoder features, enforcing internal consistency
rather than external style. The layer acts at $16\times16$ and $32\times32$ because that
is where iris colour is decided; a global vector at the $4\times4$ bottleneck and
attention there were both tried and did not move the defect (\S\ref{sec:arch}).

\emph{History of the definition.} The first implementation updated $h$ sequentially over
regions (later regions pooled features already offset by earlier ones; regions overlap,
as the eyes lie inside the face hull). The parallel form \eqref{eq:rp} differs from it
by a mean $|\Delta\xh|$ of $0.008$ ($1.6\%$ of $|\xh|$) on a trained checkpoint; the
positive result was re-established on the parallel form before the final run
(\S\ref{sec:condres}).

\subsection{Layout prior}\label{sec:prior}

\begin{figure}[t]
\centering
\diagram{pipeline}{\linewidth}
\caption{The system. (A) The generator: flow matching with x-prediction on the
training images and their detected layouts, then distillation of its midpoint-8
sampler into a one/two-jump flow map. (B) The layout prior, fitted independently
from the same detected landmarks: a point-distribution model separates pose from
shape, PCA keeps the 99\% shape subspace, and a Gaussian mixture is fitted over
pose and PCA coefficients. (C) Inference draws a layout from the prior (optionally
edited), rasterises it to masks and renders in one or two network evaluations; no
real image is involved.}
\label{fig:pipeline}
\end{figure}

Unconditional generation requires a layout that comes from no real image. The prior
(\texttt{data/layouts.py}, fitted by \texttt{experiments/landmark\_prior.py}) is a
point-distribution model \citep{cootes1995asm} with a Gaussian-mixture density.

\emph{Parametrisation.} For a landmark set $L\in\Rb^{28\times2}$: pose
$(c_x,c_y,\log s,\vartheta)$ with $c$ the centroid, $s$ the inter-ocular distance
(between the two eye centroids) and $\vartheta$ the angle of the eye line; shape $=$
the points centred, rotated by $-\vartheta$ and divided by $s$, flattened to $\Rb^{56}$.
The map is exactly invertible (round-trip error $2\times10^{-7}$).

\emph{Mirror augmentation.} $x\to-x$ with the point permutation that reverses every
horizontal run and swaps the left/right groups
($[4,3,2,1,0,\ 10,9,8,\ 7,6,5,\ 19,18,17,22,21,20,\ 13,12,11,16,15,14,\ 23,\ 26,25,24,27]$).
A wrong pairing shows as a large residual between the mean layout and the mean mirrored
layout; the verified permutation gives a maximum per-point residual of $0.085$ (the
remainder is genuine dataset asymmetry: the image-left eye is systematically smaller,
$0.356$ vs $0.383$, i.e.\ slightly turned faces).

\emph{Fit.} On the 20{,}519 layouts with mean confidence $\ge0.7$, doubled by mirroring:
PCA on the shape vectors keeping $99\%$ of variance (32 components), then a
Dirichlet-process Gaussian mixture (variational, up to 16 components, full covariances,
500 iterations; \texttt{sklearn.BayesianGaussianMixture}) on $[\text{pose};\ \text{PCA
coefficients}]$; 16 components remain active. Sampling needs only numpy: choose a
component by weight, draw from its Gaussian, invert the PCA and the pose normalisation,
rasterise as in \S\ref{sec:masks} (out-of-canvas points clipped only at rasterisation).

\emph{Checks} (4{,}000 samples each; Table~\ref{tab:prior}): the sampled marginals match
the real ones, including the left/right openness \emph{difference}, and a
jitter-the-real-layouts baseline gives the same picture. Grids of sampled masks are
plausible by eye, including the closed-eye mode.

\begin{table}[t]
\centering\small
\caption{Layout prior: marginal statistics (mean / std, $[-1,1]$ units) of real
layouts, DP-GMM samples and a KDE baseline (real layouts + Gaussian jitter,
$\sigma=0.01$).}
\label{tab:prior}
\begin{tabular}{lccc}
\toprule
measure & real & DP-GMM & KDE \\
\midrule
eye spacing & 0.769 / 0.055 & 0.768 / 0.054 & 0.769 / 0.057 \\
eye size A / B & 0.356 / 0.383 & 0.367 / 0.369 & 0.369 / 0.369 \\
openness A / B & 0.604 / 0.564 & 0.583 / 0.582 & 0.587 / 0.588 \\
$|$openness A $-$ B$|$ & 0.137 / 0.131 & 0.131 / 0.122 & 0.148 / 0.140 \\
mouth width & 0.185 / 0.115 & 0.185 / 0.108 & 0.182 / 0.113 \\
face width & 1.424 / 0.092 & 1.422 / 0.091 & 1.422 / 0.092 \\
off-canvas fraction & 0.003 & 0.001 & 0.004 \\
\bottomrule
\end{tabular}
\end{table}

\subsection{Optimisation}\label{sec:optim}

Adam with global-norm gradient clipping at $1.0$; batch 128; learning rate $10^{-3}$
peak with a linear warm-up over 500 steps (clamped to $10\%$ of short runs) and a cosine
decay to $10^{-5}$ at the last step (\texttt{optax.warmup\_cosine\_decay\_schedule}).
Weights published for sampling and evaluation are an exponential moving average with
decay $0.999$, warmed up as $\min(0.999,(1+s)/(10+s))$ at step $s$ so early averages
are not dominated by the initialisation (the scheme of the DDPM reference code,
\citealp{ho2020ddpm}; the decay is shorter than the literature's $0.9999$ because these
runs are tens of thousands of steps, not hundreds of thousands). Seed 49 for every
run; the dataset order, noise and $t$ draws are all derived from it.

The schedule is a late addition: every earlier run used a \emph{constant} $10^{-3}$,
which turned out to be at the edge of stability for the $37$\,M model. Three runs with an
added module diverged in a single epoch after tens of epochs of steady descent
(\S\ref{sec:arch}, \S\ref{sec:final}); the plain model survived 200 epochs with a thin
margin. Constant $3\times10^{-4}$ gives a \emph{lower} training loss but a \emph{worse}
FID at 30 epochs ($63.9$ vs $56.5$), so a lower constant rate was not adopted;
warm-up plus decay was.

% =====================================================================
\section{Evaluation protocol}\label{sec:protocol}

\paragraph{FID.} Inception pool features (\texttt{fidax}), 5{,}000 generated images
against statistics of the full 21{,}551-image set. Real-vs-real at 5{,}000 samples gives
a floor of $2.1$. \textbf{Noise:} two seeds of the same 30-epoch configuration scored
$56.5$ and $52.7$; differences under $\approx5$ FID at short horizons are not evidence.
\textbf{Sampler:} on the 200-epoch model, 16 / 32 / 64 Dopri5 steps score
$31.6$ / $31.4$ / $32.0$; training-time and screening FID therefore use 16 steps at a
quarter of the cost. Training batch 256 is \emph{slower} per epoch than 128
($39.9$ vs $29.0$\,s) on this GPU; 128 is kept.

\paragraph{Loss as the sensitive statistic.} Across seeds the 30-epoch training loss
varies by $\pm0.0003$ while FID varies by $\pm4$; for architecture changes that leave
the objective unchanged, loss is the sharper screen.

\paragraph{Precision / recall.} \citet{kynkaanniemi2019pr}, $k=3$ nearest neighbours in
Inception space, 5{,}000 vs 5{,}000; real-vs-real ceiling $0.774$ / $0.765$.

\paragraph{Iris-mismatch rate.}\label{sec:eyemetric} Split the dataset-mean eye mask at
the vertical midline into a left and a right region. For an image, take the
mask-weighted mean chroma $(C_b,C_r)$ (BT.601, from $[-1,1]$ RGB) of each region and
their Euclidean distance. A sample is ``mismatched'' if that distance exceeds the real
data's 95th percentile ($0.1027$), so real data scores $\approx5\%$ by construction
(two disjoint halves: $5.1\%$ and $5.9\%$). Validated by inspection: the 32 most-flagged
samples of the 200-epoch model are all genuine red/blue, yellow/blue, purple/green pairs
in otherwise well-formed faces; the least flagged are matched. \textbf{This metric is
uncorrelated with FID}: the best-FID unconditioned model has a $35.5\%$ rate.

\paragraph{Local-detail proxies.} Mean Sobel magnitude (``edge mass'') and mean ink map
of samples relative to real images; and, for the bottleneck analysis, per-region
prediction error and per-cell texture statistics (\S\ref{sec:bottleneck}).

% =====================================================================
\section{Auxiliary losses on $\xh$: theory}\label{sec:theory}

The shipped model added to \eqref{eq:loss} an edge loss and (optionally) a line loss:
\begin{align}
\mathcal L_{\text{sobel}}&=\tfrac12\bigl(\E|\nabla_x\xh-\nabla_x x_1|+\E|\nabla_y\xh-\nabla_y x_1|\bigr),
\quad\text{Sobel gradients per channel, L1 \citep{mathieu2016gdl}},\\
\mathcal L_{\text{line}}&=\E\,|\mathrm{ink}(\xh)-\mathrm{ink}(x_1)|,\qquad
\mathrm{ink}(x)=\mathrm{relu}\bigl(\mathrm{softclose}_{3\times3}(\mathrm{luma}\,x)-\mathrm{luma}\,x\bigr),\ \beta=10,
\end{align}
with weights $0.1$ and $0$ respectively.

\begin{proposition}[Irreducibility]
Let $x^*(t)=\E[x_1\mid x_t]$ be the minimiser of \eqref{eq:loss}. For any
$\mathcal L_{\text{aux}}(\xh,x_1)\ge0$ with $\mathcal L_{\text{aux}}(x_1,x_1)=0$, the
value $\mathcal L_{\text{aux}}(x^*(t),x_1)$ is in general strictly positive and cannot
be reduced by any network without leaving the optimum of \eqref{eq:loss}. Only the
excess $\mathcal L_{\text{aux}}(\xh)-\mathcal L_{\text{aux}}(x^*)$ is a usable signal.
\end{proposition}

At low $t$, $x^*$ is a blurry average; an edge or line loss against the sharp $x_1$
measures that blur, and its gradient at $x^*$ points toward ``sharper than the evidence
supports''. Whether this is \emph{bias} rather than harmless noise depends on the
minimiser of the auxiliary term: L1 on gradient differences is minimised by the
per-pixel \emph{median} of the target gradient, which under uncertainty about edge
position is nearer zero than the mean, so minimising an irreducible L1 penalty rewards
\emph{flatter} predictions (hedging). The line loss is L1 on a nonlinear map and is
biased for that reason too. Section~\ref{sec:offline} estimates the reducible fraction
empirically; \S\ref{sec:ablations} confirms the prediction in training.

\begin{proposition}[Weighted losses]\label{prop:weight}
For a per-pixel weight $w$ and prediction $c$ at fixed $x_t$,
\[
J(c)=\E[w\,(c-Y)^2\mid x_t],\qquad \nabla J=0\ \Rightarrow\
c^*=\frac{\E[wY\mid x_t]}{\E[w\mid x_t]}=\E[Y\mid x_t]+\frac{\mathrm{Cov}(w,Y\mid x_t)}{\E[w\mid x_t]} .
\]
$c^*=\E[Y\mid x_t]$ iff $w$ is conditionally independent of $Y$ given $x_t$, i.e.\ a
function of $(t,x_t)$ only.
\end{proposition}

\emph{Proof sketch.} Tower property, condition on the argument of the predictor (not on
$Y$: conditioning on $Y$ would ``derive'' $c=Y$, which no function of $x_t$ can be);
$J$ is a strictly convex quadratic in $c$; differentiate and expand
$\E[wY]=\E[w]\E[Y]+\mathrm{Cov}(w,Y)$. Equivalently, conditioning on the pair
$(x_0,x_1)$: a weight $w(x_0,x_1)$ is the unweighted CFM loss under the re-weighted
coupling $\tilde p\propto w\,p$, so by \citet{lipman2023fm} (Thm.~2) the learned field
transports noise to the data distribution \emph{tilted by $w$}. The per-sample gradient
factorisation $W\odot r$ is correct but does not make the original optimum stationary,
because $\E[w\,r\mid x_t]=\mathrm{Cov}(w,r\mid x_t)\neq0$ for a stochastic coupling
($r=0$ for every sample is unattainable). Weights of the form $\lambda(t)$---the
$1/(1-t)^2$ factor, any $t$-gate---are free; weights built from $x_1$ (Sobel magnitude,
ink map, region masks) are not, unless the map is nearly constant across the posterior
(true for the eye mask on aligned faces, false for edge maps at low $t$).

% =====================================================================
\section{Offline tests (no training)}\label{sec:offline}

All on the 100-epoch $9.4$\,M checkpoint that was trained \emph{with} the edge loss
(a caveat noted below), 256 images, seed 0 (\texttt{experiments/aux\_signal.py} on the
archive branch).

\paragraph{Discrimination.} $\mathcal L(x,d(x))/\mathcal L(x,x')$ for degradations $d$
and a random other image $x'$ (Table~\ref{tab:disc}). Both losses are $5$--$6\times$ more
sensitive than pixel L2 to stroke erasure (a $3\times3$ luma closing) and to a 1-px
jitter, and blind to global colour: they measure what they claim.

\begin{table}[t]
\centering\small
\caption{Discrimination: loss under a degradation, normalised by the loss against an
unrelated image.}
\label{tab:disc}
\begin{tabular}{lccc}
\toprule
degradation & pixel L2 & sobel & line \\
\midrule
blur $\sigma=1$ & 0.064 & 0.379 & 0.552 \\
noise $\sigma=0.05$ & 0.005 & 0.112 & 0.103 \\
desaturate 50\% & 0.019 & 0.075 & 0.000 \\
brightness $+0.2$ & 0.073 & 0.000 & 0.000 \\
erase strokes & 0.097 & 0.355 & 0.550 \\
jitter 1 px & 0.206 & 0.625 & 0.812 \\
\bottomrule
\end{tabular}
\end{table}

\paragraph{Pixel-space descent.} From $\text{blur}_1(x)$, 200 Adam steps on the pixels
minimising $\text{L2}+w\,\mathcal L_{\text{aux}}$: L2 alone recovers ink mass by step 25
and $113$\,dB PSNR by step 200; every auxiliary term is slower (sobel $0.1$: $60$\,dB;
line $1.0$: $47$\,dB). With a known target the terms carry nothing L2 lacks.

\paragraph{Irreducible fraction.} For each $t$ on a $0.05$ grid, match the checkpoint's
$\xh(t)$ with a Gaussian blur of $x_1$ at equal PSNR and read the blur's auxiliary loss
as an estimate of $\mathcal L_{\text{aux}}(x^*)$ (Table~\ref{tab:gate}). The line loss of
$\xh$ is \emph{below} the matched blur at every $t$; the sobel loss is within $\pm5\%$
except a $5$--$10\%$ excess for $t\in[0.6,0.85]$. At every $t$ at least $90\%$ of both
penalties is unreachable, and no $t$-window contains signal worth gating for. With the
default weights the auxiliary penalty is $0.6$--$1.1\times$ the velocity loss for
$t<0.5$. Caveats: a Gaussian blur is not $\E[x_1\mid x_t]$ (it likely \emph{over}-states
the irreducible part), and the checkpoint was trained with the edge loss; hence the
ablations of \S\ref{sec:ablations}.

\begin{table}[t]
\centering\small
\caption{Irreducibility estimate: auxiliary loss of $\xh(t)$ versus that of the
PSNR-matched blur $\text{blur}_{\sigma^*}(x_1)$, and the reducible fraction.}
\label{tab:gate}
\begin{tabular}{lcccc}
\toprule
$t$ & $\sigma^*$ & sobel: $\xh$ / blur & line: $\xh$ / blur & reducible sobel / line \\
\midrule
0.10 & 8.0 & 0.807 / 0.820 & 0.099 / 0.101 & 0 / 0 \\
0.30 & 2.5 & 0.677 / 0.744 & 0.090 / 0.098 & 0 / 0 \\
0.50 & 1.1 & 0.496 / 0.516 & 0.071 / 0.085 & 0 / 0 \\
0.65 & 0.7 & 0.367 / 0.344 & 0.054 / 0.068 & 0.06 / 0 \\
0.85 & 0.5 & 0.205 / 0.184 & 0.029 / 0.041 & 0.10 / 0 \\
\bottomrule
\end{tabular}
\end{table}

\paragraph{Gradient geometry.} Cosine similarity of per-sample gradients with respect
to $\xh$: $\cos(\nabla\text{L2},\nabla\text{sobel})\approx0.65$,
$\cos(\nabla\text{L2},\nabla\text{line})\approx0.40$ at every point tested. The sobel-L1
gradient is $\text{Sobel}^\top\mathrm{sign}(\cdot)$: a saturated $\pm1$ stripe pattern
over the whole image including flat background.

% =====================================================================
\section{Training ablations of the losses}\label{sec:ablations}

$9.4$\,M model, seed 49, 40 epochs, FID every 10 unless stated.

\paragraph{Shipped default vs off (paired).} Table~\ref{tab:edge}: turning the edge loss
\emph{off} is better at every evaluation and the model trained \emph{with} it generates
\emph{fewer} edges (edge mass $0.76$ vs $0.86$ of real; ink $0.68$ vs $0.79$). On
training-time interpolants the two models agree within $1.5\%$
($\mathcal L_{\text{sobel}}(\xh,x_1)$ at $t=0.1$: $0.808$ vs $0.821$)---the sliver
\S\ref{sec:offline} called reducible---so the difference is how hedged predictions
compound along the sampling trajectory. Decision: default $0$.

\begin{table}[t]
\centering\small
\caption{Loss ablations, $9.4$\,M model, 40 epochs. ``Class'' is the verdict of
Prop.~\ref{prop:weight}.}
\label{tab:edge}
\begin{tabular}{lcccl}
\toprule
arm & FID@39 & ink/real & edge/real & class \\
\midrule
no auxiliary (reference) & 68.0 & 0.79 & 0.86 & --- \\
sobel L1, weight 0.1 (shipped default) & 74.7 & 0.68 & 0.76 & biased (median) \\
sobel L2, weight 0.05 & 68.0 & -- & -- & unbiased (linear map) -- inert \\
FM weight $1+|\text{Sobel}(x_1)|/\text{mean}$ & 65.3 & 0.89 & 0.91 & biased (tilt) \\
FM weight from $|\text{Sobel}(\text{sg}\,\xh)|$ & 65.3 & 0.79 & 0.87 & unbiased \\
FM weight $1+\text{ink}(x_1)/\text{mean}$ & 97.0 & & & biased; $\sim50\times$ per pixel \\
FM weight $1+m_{\text{eye}}$, coef.\ 1.0 / 0.2 & 117.7 / 76.0 & & & $\sim$unbiased; $20\times$ / $5\times$ \\
\bottomrule
\end{tabular}
\end{table}

\paragraph{Variants.} The L2 Sobel term is exactly inert (68.0 vs 68.0), as predicted
for an unbiased term with $\le10\%$ reducible signal. The two edge re-weightings end at
the same FID whether or not they carry the covariance term of Prop.~\ref{prop:weight};
the tilt shows only in the proxies (edgier samples). Sparse-map weights (ink, eye) are
harmful at every coefficient tried because the batch-mean normalisation turns
``average extra weight 1'' into $20$--$50\times$ on a few percent of pixels; the eye
weight does reduce iris mismatch dose-dependently ($41.8\to37.5\to32.8\%$) at a
prohibitive FID cost.

\paragraph{At length.} $37$\,M model, 200 epochs, FID at 49/99/149/199: plain
$43.1/36.3/36.5/32.0$; with the unbiased edge re-weighting $42.9/37.7/35.9/33.4$. The one
loss-side candidate is a wash. A mask-prediction auxiliary head on the encoder (BCE on
area-pooled masks, weights $0.25$ and $1.0$) was monotonically worse at $9.4$\,M over 100
epochs ($57.6$, $69.7$ vs $53.5$ best) and at $37$\,M ($62.4$ vs $56.5$). \textbf{The
pixel-space auxiliary family is closed} at both widths and all horizons tested.

% =====================================================================
\section{Scaling and the noise floor}\label{sec:scaling}

Width and length are the two levers that moved FID substantially: $9.4$\,M / 40 epochs
$68.0$; $9.4$\,M / 100 epochs $53.5$; $37$\,M / 30 epochs $56.5$; $37$\,M / 200 epochs
$32.0$, still falling ($36.5\to32.0$ over the last 50). A spectral check supports
spending on width: the effective rank (entropy of the normalised singular values of each
convolution reshaped to $(\text{out},\text{in}\cdot k^2)$) is $0.88$ of the maximum at
$9.4$\,M and $0.85$ at $37$\,M, i.e.\ the doubled width is used, not idle; from 30 to
200 epochs the stable rank rises ($13.8\to17.5$) while the entropy rank falls slightly
($0.853\to0.825$), the signature of training concentrating a few directions while
spreading energy into the mid-spectrum. Curation (\S\ref{sec:data}) gave $53.2$ / loss
$0.0875$ against $56.5$--$52.7$ / $0.0890$ at 30 epochs.

% =====================================================================
\section{Sampling-time edge guidance}\label{sec:guidance}

Priors about \emph{samples} belong at sampling time, where they cannot bias the learned
field. With $E(\xh)=\mathrm{relu}\bigl(m_{\text{real}}-\overline{|\text{Sobel}(\xh)|}\bigr)^2$
(the one-sided deficit against the data's mean edge mass, $m_{\text{real}}=1.30$),
\begin{equation}
g=\frac{\partial E(\xh(x_t,t))}{\partial x_t},\qquad
\xh'=\xh-\lambda\,\mathbb 1[t_{\text{lo}}<t<t_{\text{hi}}]\,\frac{g}{\mathrm{rms}(g)},\qquad
v=\frac{\xh'-x_t}{\max(1-t,\epsilon_d)},
\end{equation}
gradient through the network, RMS-normalised per sample. Results (5{,}000-sample FID):
$9.4$\,M model $67.4\to59.0$ at $\lambda=0.02$ ($\lambda=0.05/0.10/0.20$:
$59.4/62.1/69.5$, over-shooting as edge mass passes real); $37$\,M 200-epoch model
$31.4\to28.1$. Restricting guidance to a $t$-interval \citep{kynkaanniemi2024interval}
helps slightly: on the $9.4$\,M model $\lambda=0.05$ on $t\in(0.5,1)$ gives $57.8$,
while $(0,0.5)$ is the worst window ($61.4$). Guidance raises precision
($0.52\to0.58$) and leaves recall at $0.15$: a mode-seeking correction, which bounds its
value at a few FID. Cost: one backward pass per network evaluation. On the
region-pool model of \S\ref{sec:final} (\texttt{best\_model.eqx}, $34.6$) the gain is
the same: $32.5$ at $\lambda=0.02$ on the full window, $32.8$--$33.6$ for the other
$(\lambda,\text{window})$ pairs; precision $0.47\to0.49$, recall $0.125\to0.108$.

\subsection{Autoguidance}\label{sec:autoguide}

\citet{karras2024autoguidance} guide a model with a worse version of itself,
\begin{equation}
v = v_{\text{good}} + w\,(v_{\text{good}} - v_{\text{bad}}),
\end{equation}
where $v_{\text{bad}}$ is a smaller or less-trained model of the same family: the
extrapolation removes the errors the two models share and the bad one makes more
strongly. For a layout-conditioned model both velocities see the same mask; the cost is
one extra forward pass per evaluation and no training
(\texttt{experiments/autoguide.py}). Good model: \texttt{wide\_rp\_300/best\_model.eqx}
($34.6$ unguided), 5{,}000 samples, 16 steps, prior layouts.

\begin{center}\small
\begin{tabular}{llcccccc}
\toprule
bad model & window & $w{=}0.5$ & $0.75$ & $1$ & $1.25$ & $1.5$ & $2$ \\
\midrule
\texttt{confirm\_rp} (same recipe, 40 ep, FID 52) & $(0,1)$ & 24.9 & 22.0 & \textbf{21.2} & 23.3 & 29.7 & 56.9 \\
\texttt{confirm\_rp} & $(0.3,1)$ & 26.6 & & 23.1 & & & 44.7 \\
epoch 49 of the same run & $(0,1)$ & 28.9 & & 26.5 & & & 51.6 \\
\bottomrule
\end{tabular}
\end{center}

At the optimum ($w=1$, full window): 64 sampler steps give the same $21.2$; real
layouts $20.8$ (vs $34.0$ unguided, so the prior-vs-real gap is unchanged); edge mass
rises from $0.90$ to $1.08$ of real; on the best checkpoint of the one-segment run
(\S\ref{sec:final}, $33.2$ unguided) the result is $\mathbf{18.1}$. Three readings.
(i) $-13$ FID is three to four times the gain of any energy guidance and moves the
layout-conditioned model from $3.5$ FID behind the schedule-matched unconditioned control
to $11$ ahead of it. (ii) The guide must be bad in the right way: a separately trained
40-epoch checkpoint of the same recipe beats epoch 49 of the \emph{same} run---which
shares the good model's initialisation and early trajectory---by $5$ FID at every $w$.
(iii) The response is sharp above $w=1$ (collapse at $2$), and restricting the window,
which helped energy guidance, hurts here. The rise in edge mass says what is being
corrected: the hedged, low-contrast detail that \S\ref{sec:bottleneck} identifies as the
model's characteristic error, which an under-trained model makes more strongly. Combining
with one refinement round (\S\ref{sec:refine}) \emph{hurts} ($25.2$). Iris mismatch,
precision/recall and mask-following under autoguidance have not yet been re-measured.
The result reproduces on a second dataset (\S\ref{sec:celeba}).

% =====================================================================
\section{Where the model fails}\label{sec:bottleneck}

On the $37$\,M 200-epoch unconditioned model (\texttt{experiments/bottleneck.py}).

\paragraph{Kind of gap.} FID floor $2.1$. Precision $0.521$, recall $\mathbf{0.154}$
(ceiling $0.774$/$0.765$). Coverage dominates: $85\%$ of real images have no generated
neighbour. The samples farthest from the real manifold are washed-out, low-contrast
images with incoherent hair, clothing and background---averages of modes the model has
not learned, not broken faces.

\paragraph{Where.} Prediction error divided by per-pixel data variance (512 real pairs):
at $t=0.3$, face $0.41$, eyes $0.43$, \textbf{mouth $0.50$}, hair/background $0.25$; the
mouth is recovered worst relative to what is predictable because small low-variance
high-frequency features receive the least loss weight. Generated-vs-real local
statistics on an $8\times8$ grid: edge mass $-9$ to $-12\%$, luma std $-9$ to $-13\%$
(worst at the mouth/nose centre, $-30$ to $-40\%$), chroma std $-16$ to $-21\%$
everywhere. Colour \emph{means} are right (skin chroma $-0.062/0.089$ vs
$-0.070/0.098$); colour \emph{spread} is $13\%$ narrow and the chroma-histogram distance
to real is $2.1\times$ the real-vs-real floor. A palette loss therefore has no target;
what is missing is diversity, which no per-sample loss supplies.

\paragraph{Eyes.} Iris mismatch $35.5\%$ on this model (\S\ref{sec:eyemetric}), i.e.\ a
long-range colour-coupling failure that FID does not see and that motivated
\S\ref{sec:layout}.

\paragraph{Region-pool model.} The same analysis on \texttt{wide\_rp\_300/best\_model.eqx}
(\S\ref{sec:final}): precision $0.471$, recall $0.125$---the same coverage picture,
slightly worse on both axes at a similar FID. With the layout given, the
variance-normalised error at $t=0.3$ is even across regions (face $0.32$, eyes $0.30$,
mouth $0.33$, hair $0.23$; the mouth is no longer the outlier), and the local-statistics
deficit is a uniform $-9$ to $-12\%$ on edge mass, luma and chroma std: the layout
removes the \emph{where} uncertainty and leaves the hedging.

% =====================================================================
\section{Architecture variants that did not help}\label{sec:arch}

All paired, $37$\,M, 30 epochs unless stated; loss (stable to $\pm0.0003$) is the
primary read-out.

\begin{itemize}[leftmargin=1.5em,itemsep=1pt]
\item \textbf{Standard skips.} $9.4$\,M: $76.0$ vs $76.0$ (neutral). $37$\,M at
      constant $10^{-3}$: FID $188$, loss flat at $0.108$ from epoch 3, samples frozen
      into a tiled texture---an optimisation collapse; at $3\times10^{-4}$: $63.6$ with
      loss $0.0874<0.0890$ but no better than the legacy layout. Kept legacy.
\item \textbf{Mid-block (ResBlock--attention--ResBlock at $4\times4$).} $9.4$\,M, 40
      epochs: $78.7$ vs $77.1$. $37$\,M, 200 epochs, constant $10^{-3}$: $75.8$ with a
      loss spike to $0.635$ at epoch 12---an optimisation failure, not a verdict on
      attention. It had the lowest iris mismatch of the unconditioned models ($23.0\%$),
      consistent with global mixing helping the coupling.
\item \textbf{Global code.} $g=W_2\,\mathrm{SiLU}(W_1\,\mathrm{LN}(\text{mean}_{hw}z)+b_1)+b_2$
      from the bottleneck $z$, $W_2=0$ at init, added to the decoder's time embedding
      (the ADM/DiT class-embedding recipe). Loss $0.0869$ vs $0.0870$, FID $63.8$ vs
      $63.9$, mismatch $43.8\%$ vs $42.6\%$. The path \emph{did} switch on
      ($\|g\|\approx0.6$--$0.8\,\|e\|$, image-dependent) and changed nothing: the
      decoder is not short of global \emph{information}; the defect is in how two
      regions are \emph{rendered}, which is what \S\ref{sec:regionpool} addresses.
\end{itemize}

% =====================================================================
\section{Layout conditioning and region pooling: results}\label{sec:layout}

\subsection{Conditioning alone}\label{sec:condres}

$37$\,M, curated. With dropout $0.15$ (30 epochs): unconditional mode $62.6$, prior
layouts $57.8$, real layouts $57.2$, mismatch $29.7\%$; without dropout (40 epochs, loss
$0.0813$): prior $50.7$, real $50.1$, mismatch $33.6\%$. Unconditional references at
30 epochs: $56.5/52.7/53.2$. Layout information lowers the training loss but not FID
(real and prior layouts agree, so the prior is not the limitation) and does not touch
iris agreement: a layout says where the eyes are, not what colour they share. Dropout
costs $\approx7$ FID at this horizon.

\subsection{Region pooling}

Same recipe (pure conditional, 40 epochs, LR $10^{-3}$ constant, seed 49, 16-step FID):

\begin{center}\small
\begin{tabular}{lccc}
\toprule
40 epochs & loss & FID prior / real layouts & iris mismatch prior / real \\
\midrule
no region pool & 0.0813 & 50.7 / 50.1 & 33.6\% / 38.7\% \\
region pool (sequential form) & 0.0807 & 48.2 / 48.0 & \textbf{6.6\%} / 5.1\% \\
region pool (parallel form, confirmation) & 0.0815 & 52.1 / 51.4 & \textbf{6.25\%} / 5.1\% \\
real data & & & 5.0--5.9\% \\
\bottomrule
\end{tabular}
\end{center}

The mismatch rate falls to the data's own rate. Paired samples (same noise, same
layouts) are the same face pair-for-pair with only the iris pairs changed---matched and
more saturated. The FID difference between the two region-pool runs ($48.2$ vs $52.1$)
is inside the single-seed noise. The eye-region chroma distribution of the region-pool
model matches real data in spread (saturation std $0.067$ vs $0.071$; $C_b$, $C_r$ std
$85$--$95\%$ of real) and in hue histogram (same dominant blue-green band, $32\%$ in
both), so the layer enforces \emph{agreement}, not a particular colour.

\paragraph{Mask-following.} Does the model draw the face where the layout says?
512 prior layouts are rendered by each model, the landmark detector is re-run on the
samples, and the detected landmarks are scored against the layout the image was
rendered from (\texttt{experiments/mask\_following.py}; real images against their own
landmarks are the ceiling).

\begin{center}\small
\begin{tabular}{lcccc}
\toprule
model & landmark error all / eyes / nose / mouth & det.\ conf.\ nose / mouth & IoU face / eyes \\
\midrule
conditional, no region pool (40 ep) & 0.052 / 0.027 / 0.101 / 0.062 & 0.57 / 0.71 & 0.918 / 0.830 \\
conditional + region pool (\texttt{best\_model}) & \textbf{0.046 / 0.026 / 0.091 / 0.053} & \textbf{0.64 / 0.75} & \textbf{0.931 / 0.850} \\
real & 0 & 0.82 / 0.85 & 1 \\
\bottomrule
\end{tabular}
\end{center}

Both models follow the layout closely (100\% detection; the eye error of $0.027$ in
$[-1,1]$ units is under one pixel at $64$\,px). Region pooling tightens nose and mouth
placement by $\approx10\%$ and raises the detector's confidence in those features, but
the remaining gap to real in nose/mouth legibility ($0.64$ vs $0.82$) is the larger
effect and is not a layout problem.

% =====================================================================
\section{Final training run}\label{sec:final}

Recipe (\texttt{just paper}): $37$\,M U-Net with region pooling, pure conditional on the
three-channel layout, curated data, batch 128, EMA $0.999$, Dopri5 16-step FID every 50
epochs with prior-sampled layouts, 300 epochs in total.

\paragraph{Course of the run.} The first attempt ran at a constant $10^{-3}$ and
diverged in a single epoch at epoch 80 (loss $0.0784\to0.2837$, samples destroyed,
partial recovery to $0.107$), after a steady descent to $0.0776$ and FID $42.0$ at epoch
49. The warm-up-plus-cosine schedule of \S\ref{sec:optim} was introduced in response and
the run resumed from the epoch-49 EMA weights with a fresh optimiser (Adam moments are
not checkpointed). A second segment of 49 epochs (loss $0.0770$) ended in a
\texttt{cuSolver} failure inside the FID computation caused by a concurrent GPU job---an
operational error, not a modelling one---and the run resumed again from those weights
for the final 200 epochs. The effective schedule is three segments (49 constant, 49 and
200 warm-up-cosine), with the learning rate near its peak at both cut points, so it
approximates one long cosine; checkpoints and loss logs of every segment are kept
(\texttt{runs/wide\_rp\_300\_ep49.eqx}, \texttt{\_ep98.eqx}, and the corresponding
\texttt{losses.json} files).

\paragraph{Results.} Final loss $0.0557$. Training-time FID over the last
segment (16-step sampler, prior layouts), at cumulative epochs $\approx147/197/247/297$:
$47.8/34.9/36.2/37.5$---best at $\approx200$ cumulative epochs, drifting up
thereafter; \texttt{best\_model.eqx} is that checkpoint. Final evaluation of the last
weights (5{,}000 samples, 16 steps): FID $\mathbf{37.8}$ with prior layouts and
$34.2$ with real layouts; \textbf{iris mismatch $4.7\%$} under either (real data
$5.0\%$; mean left/right chroma distance $0.039$ vs $0.037$ real). Layout adherence
is strong (Fig.~\ref{fig:showcase}, layout-to-image panel: tilt, eye spacing and eye
height track the mask in every sample). Against the plain unconditioned 200-epoch model
($31.4$ at the same sampler) the conditioned model is $\approx6$ FID behind at the
end and $\approx3.5$ behind at its best point, while solving the eye defect outright;
the prior-vs-real layout gap ($3.5$ FID), absent at 40 epochs ($50.7$ vs $50.1$),
opened as adherence sharpened, which points at a residual mismatch between the prior's
layouts and the data's as one component of the deficit. Off-the-shelf
anime super-resolution models (Real-ESRGAN variants, APISR) were tried on the samples
for presentation and dropped: trained on clean-then-degraded pairs, they render the
generator's own errors as confident detail, and the native $64\times64$ samples read
better.

\paragraph{The two confounds, resolved.} Two further 200-epoch runs, each one segment
of warm-up plus cosine with FID every 25 epochs: \texttt{ctrl\_sched}, the plain
unconditioned model on the same schedule (the missing control), and
\texttt{wide\_rp\_200}, the region-pool recipe as recommended above.

\begin{center}\small
\begin{tabular}{lcccccccc|c}
\toprule
epoch & 24 & 49 & 74 & 99 & 124 & 149 & 174 & 199 & final loss \\
\midrule
plain, same schedule (\texttt{ctrl\_sched}) & 64.1 & 38.2 & 30.6 & 30.8 & 29.1 & 29.5 & 29.5 & \textbf{29.3} & 0.0789 \\
conditional + region pool (\texttt{wide\_rp\_200}) & 58.9 & 37.9 & 35.0 & 33.1 & \textbf{32.9} & 34.5 & 36.5 & 37.8 & 0.0537 \\
\bottomrule
\end{tabular}
\end{center}

\texttt{wide\_rp\_200/best\_model.eqx} (epoch 124) evaluates at FID $32.9$ with prior
layouts, $31.4$ with real layouts, iris mismatch $4.3\%$ under either. Three
conclusions. The schedule alone is worth $\approx2.7$ FID on the plain model
($32.0\to29.3$), whose curve is flat from epoch 75. The one-segment run reproduces the
three-segment one (best $32.9$ vs $34.9$, final $37.8$ vs $37.8$): the warm restarts
were not the cause of the drift. The drift after epoch $\approx125$ is real and specific
to the conditioned model---its loss keeps falling ($0.065\to0.054$) while the control's
loss and FID are both flat---so the layout-conditioned region-pool model stands $3.5$
FID behind the matched plain model at its best point and $8.5$ behind at the end, at
$4.3\%$ iris mismatch against $35\%$. The drift is consistent with the model's growing
reliance on the region-pool broadcast (its contribution to $\xh$ triples over training),
which renders the prior's imperfect layouts more literally: the prior-vs-real gap is
$1.5$ FID at the best checkpoint and $3.5$ at the end. Autoguidance
(\S\ref{sec:autoguide}) more than recovers the deficit: $18.1$ on this checkpoint.

\subsection{Distilled flow map: one or two evaluations}\label{sec:distill}

\begin{figure}[t]
\centering
\diagram{distill}{\linewidth}
\caption{Fixed-teacher distillation. The velocity model under the midpoint
sampler is a deterministic map from (noise, layout) to an image; the student -- the
same U-Net with a second time input -- regresses its jumps $0\to1$, $0\to\tfrac12$ and
$\tfrac12\to1$ onto the teacher's states with a plain L2. Self-consistency
objectives (MeanFlow, shortcut) either diverged or saturated when fine-tuned from
this model (\texttt{experiments/METHODS.md} \S3.4).}
\label{fig:distill}
\end{figure}

Sixteen Dopri5 steps cost $\approx96$ network evaluations; an explicit midpoint
scheme with 8 steps (16 evaluations) matches it (FID $27.1$ vs $29.3$ on the plain
model) and 4 steps (8 evaluations) gives $28.1$, so below 8 evaluations a learned
\emph{flow map} $x_s = x_t + (s-t)\,u_\theta(x_t,t,s)$ is needed. Distilling the
midpoint-8 sampler of the region-pool model into such a map (Fig.~\ref{fig:distill};
100{,}000 teacher trajectories, 35 epochs) gives FID $53.0$ in one evaluation and
$\mathbf{37.8}$ in two, against the teacher's $32.9$; exported to ONNX the two-jump
sampler renders an image in $99$\,ms on two CPU threads, which is what the public
demo runs.

\subsection{Re-noise-and-re-solve refinement}\label{sec:refine}

A finished sample is re-noised to $t_0$, $x_{t_0}=t_0\,x+(1-t_0)\,\epsilon$, and the ODE
is solved again from $t_0$ with the same layout (SDEdit, \citealp{meng2022sdedit}, on
the model's own output; one round of Restart sampling, \citealp{xu2023restart}).
On \texttt{best\_model.eqx} with 16 steps per pass: $t_0=0/0.5/0.7/0.85$ gives FID
$34.6/33.8/32.8/\mathbf{31.8}$ with edge mass unchanged ($0.89$--$0.90$ of real).
Unlike guidance it adds no contrast; it is a second draw of the late-decided detail.
It does not stack with autoguidance ($21.2\to25.2$ with one round at $t_0=0.85$: the
guided field is already off the model's own manifold, and refinement pulls it back).
It remains the in-distribution alternative to the external upscalers above.

\begin{figure}[t]
\centering
\figorplaceholder{showcase.png}{\linewidth}
\caption{Showcase page produced by \texttt{generate\_figures.py}: loss and FID curves,
samples with layouts from the prior, layout-to-image pairs, real vs generated, noise
interpolation at a fixed layout, and eye-region crops.}
\label{fig:showcase}
\end{figure}

% =====================================================================
\section{Second dataset: CelebAMask-HQ}\label{sec:celeba}

Does the recipe transfer to real faces? CelebAMask-HQ \citep{lee2020maskgan} provides
30{,}000 aligned CelebA-HQ faces with 19-class segmentation maps; we use the 256\,px
mirror, area-averaged to $64\times64$, and build the
four-channel layout from label unions: face (skin, nose, brows, eyes, lips, mouth,
glasses), eyes, mouth (lips and interior), nose; region pooling adds $1-\text{face}$ as
before (\texttt{data/celebamask.py}). Curation drops $1.8\%$ (duplicate thumbnails,
empty eye/mouth/nose masks, mis-aligned or greyscale images), leaving 26{,}487. No
layout prior is fitted: sampling-time layouts are the real label maps of a held-out bank
(the last 3{,}000 ids, never trained on), so every number below is in the ``real
layouts'' mode; FID statistics are the full 30{,}000 images. At $64$\,px the two eyes
together cover one $16\times16$ cell ($14\times$ smaller than anime eyes), so the iris
mechanism is barely testable at this size; a $128$\,px run (arrays built, fits the
GPU at a $0.95$ memory fraction) was not pursued, being a different cost class for a
question the anime result already answers.

\paragraph{Training.} Same architecture ($37$\,M, \texttt{cond\_channels 4},
\texttt{region\_pool 1}), batch 128. At a constant $10^{-3}$ the run diverged at epoch
25 (loss $0.043\to0.285$, FID $41\to228$), the same failure as the anime run and earlier;
with a $5\times10^{-4}$ peak, a 2{,}000-step warm-up and cosine decay, 120 epochs reach
FID $39.9/27.3/27.3$ at epochs $39/79/119$ and a final loss of $0.036$
(\texttt{just celeba}). The FID curve is flat from epoch 80: with real layouts there is
no prior mismatch to render, and no late drift.

\paragraph{Autoguidance.} Good model: the 120-epoch checkpoint ($27.4$ unguided);
bad model: the epoch-24 checkpoint of the diverged constant-LR run (FID $41$). FID at
$w=0.8/0.9/1.0$: $\mathbf{18.6}/19.1/20.3$, edge mass $1.07/1.10/1.13$ of real
($0.96$ unguided). The anime result (\S\ref{sec:autoguide}) reproduces with a different
mask source and a different guide: $-9$ FID, optimum near $w=1$, edge mass crossing real
at the optimum.

% =====================================================================
\section{Decisions and their evidence (summary)}\label{sec:decisions}

\begin{itemize}[leftmargin=1.5em,itemsep=1pt]
\item \textbf{x-prediction, velocity-space loss, floor 0.05, logit-normal $t$}: JiT
      recipe; the parametrisation is what makes losses on $\xh$ \emph{possible}, and
      \S\ref{sec:theory} is why none are used.
\item \textbf{No auxiliary losses}: Prop.~1--2, Tables~\ref{tab:gate}--\ref{tab:edge},
      the 200-epoch wash. Closed at both widths.
\item \textbf{Guidance instead}: edge guidance $-2$ to $-8$ FID at $\lambda\approx0.02$;
      refinement $-3$ at $t_0=0.85$; \textbf{autoguidance $-13$} with a 40-epoch
      checkpoint of the same recipe at $w=1$---the largest single gain measured, and
      reproduced on CelebAMask-HQ ($-9$). All reported separately from unguided numbers.
\item \textbf{GroupNorm (8 groups), zero-init adaLN, no output gate, no final activation,
      time scale $10^3$}: documented in code; the zero-init adaLN is why every added
      module (region pool, global code) can be identity at initialisation.
\item \textbf{Legacy skips kept}: neutral at $9.4$\,M, unstable at $37$\,M.
\item \textbf{Width 64, batch 128, 16-step evaluation, curation on}: \S\ref{sec:scaling},
      \S\ref{sec:protocol}.
\item \textbf{Warm-up + cosine}: four single-epoch divergences at constant $10^{-3}$
      (three anime, one CelebA); worth $2.7$ FID on the plain model by itself.
\item \textbf{Pure conditional (no dropout)}: dropout $\approx7$ FID at the screening
      horizon; the prior makes the unconditional mode unnecessary.
\item \textbf{Region pooling at $16\times16$ and $32\times32$, four regions, zero-init,
      parallel form}: \S\ref{sec:layout}; the only structural change with a reproduced,
      metric-validated win (iris mismatch $35\%\to4.3$--$4.7\%$, mask-following up), at
      $3.5$ FID against the schedule-matched control before guidance and $11$ FID ahead
      of it after autoguidance.
\item \textbf{Layout prior: PDM + PCA(99\%) + DP-GMM}: Table~\ref{tab:prior}; sampling
      is numpy-only.
\end{itemize}

\paragraph{Open.} (1) Autoguidance is under-explored: which guide (epoch, width,
data fraction) is best---a separately trained 40-epoch checkpoint beats an early
checkpoint of the same run by $5$ FID; whether the iris rate, mask-following and recall
survive $w=1$; and whether the plain model gains as much (if not, conditioning, region
pooling and autoguidance are one recipe; if so, autoguidance is orthogonal). (2) The
conditioned model's late-training drift: a better layout prior (real landmarks with
jitter) or stopping at the FID-optimal checkpoint. (3) Coverage (recall $0.15$) remains
the largest gap of the unguided model and moves with steps and capacity; EDM-style
augmentation conditioning with hair/eye hue rotation (a dataset-specific coverage prior,
semantically valid only for this domain), spatial mixing at $16\times16$ with warm-up,
a hair-colour consistency metric for the region-pool layer, and flip-equivariant
sampling. Shelved with reasons recorded: structured output heads (line/fill, palette) as
precision-side bets with a representational risk.

% =====================================================================
\section*{Reproducibility}

Repository: \texttt{tinyflow}, branch \texttt{feat/pixel-unet}; the complete working tree
of the campaign, including every dropped variant and the offline harness, is preserved
on branch \texttt{exp/aux-loss-campaign}. \texttt{experiments/METHODS.md} holds every
table of this paper with the commands that produced it; \texttt{just paper NAME 300 50}
reproduces the final run and the figures. Hardware: one NVIDIA RTX~4090 (24\,GB);
training throughput $\approx22$\,s per epoch (166 batches of 128) for the $37$\,M model;
a 16-step 5{,}000-sample FID takes $\approx2$ minutes.

\begin{thebibliography}{9}\small
\bibitem[Cootes et al.(1995)]{cootes1995asm} T.~F. Cootes, C.~J. Taylor, D.~H. Cooper,
J.~Graham. Active shape models---their training and application. \emph{CVIU}, 1995.
\bibitem[Ho et al.(2020)]{ho2020ddpm} J.~Ho, A.~Jain, P.~Abbeel. Denoising diffusion
probabilistic models. \emph{NeurIPS}, 2020.
\bibitem[Ho \& Salimans(2022)]{ho2022cfg} J.~Ho, T.~Salimans. Classifier-free diffusion
guidance. arXiv:2207.12598, 2022.
\bibitem[Karras et al.(2024)]{karras2024autoguidance} T.~Karras, M.~Aittala, T.~Kynkäänniemi,
J.~Lehtinen, T.~Aila, S.~Laine. Guiding a diffusion model with a bad version of itself.
\emph{NeurIPS}, 2024.
\bibitem[Kynkäänniemi et al.(2019)]{kynkaanniemi2019pr} T.~Kynkäänniemi, T.~Karras,
S.~Laine, J.~Lehtinen, T.~Aila. Improved precision and recall metric for assessing
generative models. \emph{NeurIPS}, 2019.
\bibitem[Kynkäänniemi et al.(2024)]{kynkaanniemi2024interval} T.~Kynkäänniemi et al.
Applying guidance in a limited interval improves sample and distribution quality in
diffusion models. \emph{NeurIPS}, 2024.
\bibitem[Lee et al.(2020)]{lee2020maskgan} C.-H. Lee, Z.~Liu, L.~Wu, P.~Luo. MaskGAN:
towards diverse and interactive facial image manipulation. \emph{CVPR}, 2020.
\bibitem[Li \& He(2026)]{li2026jit} T.~Li, K.~He. Back to basics: let denoising
generative models denoise (JiT). \emph{CVPR}, 2026.
\bibitem[Lipman et al.(2023)]{lipman2023fm} Y.~Lipman, R.~T.~Q. Chen, H.~Ben-Hamu,
M.~Nickel, M.~Le. Flow matching for generative modeling. \emph{ICLR}, 2023.
\bibitem[Meng et al.(2022)]{meng2022sdedit} C.~Meng, Y.~He, Y.~Song, J.~Song, J.~Wu,
J.-Y. Zhu, S.~Ermon. SDEdit: guided image synthesis and editing with stochastic
differential equations. \emph{ICLR}, 2022.
\bibitem[Mathieu et al.(2016)]{mathieu2016gdl} M.~Mathieu, C.~Couprie, Y.~LeCun. Deep
multi-scale video prediction beyond mean square error. \emph{ICLR}, 2016.
\bibitem[Peebles \& Xie(2023)]{peebles2023dit} W.~Peebles, S.~Xie. Scalable diffusion
models with transformers. \emph{ICCV}, 2023.
\bibitem[Wu \& He(2018)]{wu2018gn} Y.~Wu, K.~He. Group normalization. \emph{ECCV}, 2018.
\bibitem[Xu et al.(2023)]{xu2023restart} Y.~Xu, M.~Deng, X.~Cheng, Y.~Tian, Z.~Liu,
T.~Jaakkola. Restart sampling for improving generative processes. \emph{NeurIPS}, 2023.
\bibitem[Zhu et al.(2020)]{zhu2020sean} P.~Zhu, R.~Abdal, Y.~Qin, P.~Wonka. SEAN: image
synthesis with semantic region-adaptive normalization. \emph{CVPR}, 2020.
\end{thebibliography}

\end{document}
